\documentclass[12pt]{article}
\usepackage[margin=1in]{geometry}
\usepackage{setspace}
\usepackage{amsmath, amsfonts, amssymb}
\usepackage{graphicx}
\usepackage{booktabs}
\usepackage{hyperref}
\usepackage{tgpagella}

\setstretch{1.5}

\begin{document}

\begin{center}
    \Large \textbf{Confidential Editor Report}\\[1ex]
    \large \textbf{Journal:} Journal of Environmental Economics and Management\\
    \large \textbf{Paper Title:} Monitoring the Monitors: How Environmental Information Shapes Bureaucratic Incentives in China\\
    \large \textbf{Manuscript ID:} JEEM-D-26-00040
\end{center}

\section{Recommendation and Message to Editor}

\textbf{Recommendation: Major Revision}

This paper asks an important question and has substantial empirical potential. I do not recommend acceptance in the current form, mainly because the channel interpretation is not yet sufficiently identified and the theoretical structure is not fully aligned with the empirical narrative.

My highest-priority concern is the post-2013 interpretation: the paper attributes observed shifts primarily to incentive realignment via improved observability, but multiple concurrent channels are plausible (evaluation-weight change, stronger accountability, higher central support/resources, public-awareness effects, and perception/learning effects). The revision should explicitly test or bound these alternatives.

A second core concern is the model structure: economic performance is not explicitly embedded in the promotion index. In a growth-versus-environment cadre setting, this weakens structural interpretation. A revised model should include both environmental and economic performance channels, and map directly to empirical specifications.

\section{Key Confidential Concerns}

\begin{itemize}
    \item Post-period effects may combine information precision changes with broader policy-regime changes.
    \item IV exclusion and channel interpretation need stronger robustness and falsification support.
    \item Work-report-based effort measures should be validated and complemented with more direct enforcement indicators.
    \item Promotion coding and sample exclusions (dropping demoted/retired/disciplined officials) may affect inference.
    \item Pre-period negative PM2.5 estimates and transformation choices (e.g., $\log(y+1)$ for high-pollution days) need deeper treatment.
\end{itemize}

\section{Potential Full Letter to Authors}

\subsection*{Summary}

This paper studies whether monitoring and disclosure reforms changed bureaucratic incentives for environmental performance. The topic is important and the data effort is strong. The manuscript can be significantly improved by tightening channel identification and revising the model-empirics linkage.

\subsection*{Major Concerns}

\begin{itemize}
    \item \textbf{Channel identification around post-2013 changes:}
    show stronger empirical separation of observability/incentive effects from alternative channels (accountability, weight shifts, resource support, public-awareness effects, and learning).

    \item \textbf{Model revision:}
    include explicit economic performance in the promotion index alongside environmental outcomes. Clarify that the effort FOC is a net-promotion-return condition (cost equals environmental promotion gain minus growth-related promotion loss), and discuss how post-2013 weight changes complicate empirical identification.

    \item \textbf{Effort measurement and mechanisms:}
    complement text-coded policy measures with direct enforcement indicators (penalties, supervision/inspection intensity) and provide stronger support for the administrative-measures mechanism.

    \item \textbf{Outcome and specification robustness:}
    address lateral-move interpretation, sample-selection effects from dropped officials, pre-period baseline coefficient interpretation, and alternatives to $\log(\text{HPD}+1)$.
\end{itemize}

\subsection*{Additional Suggestions}

\begin{itemize}
    \item Consider local-projection impulse-response style estimates for dynamic promotion responses.
    \item Explore hierarchy heterogeneity where possible (mayor/county/province levels).
    \item Provide year-by-year coefficient trend figures.
    \item Add heterogeneity by region, age, education, and gender.
    \item Expand appendix detail on data, variable definitions, and coding rules.
\end{itemize}

\subsection*{Closing}

I view the paper as promising but not yet ready. A focused major revision along the lines above could materially improve the contribution.

\end{document}
